\documentclass[../../../uem_mei_q.tex]{subfiles}

\begin{document}
    次の空欄中\textsf{ア}，\textsf{イ}，\textsf{ウ}，\textsf{ク}に当てはまるものを選択肢の中から選び%
    その記号をマークせよ．それ以外の空欄に当てはまる0から9までの数字を解答用紙の所定の欄に%
    マークせよ．ただし，\,\myboxb{\hspace{.5\zw}カキ\hspace{.5\zw}}\,，%
    \,\myboxb{\hspace{.5\zw}コサ\hspace{.5\zw}}\:は2桁の数，\myboxb{シスセ}\:は3桁の数である．

    \vspace{.5\zw}
    $a$を正の定数として
    \begingroup
        \setlength{\abovedisplayskip}{4pt}
        \setlength{\belowdisplayskip}{4pt}
        \setlength{\jot}{0pt}
        \begin{align*}
            &f(x)=(ax+4)^2 \\
            &g(x)=(a+4)(ax^2+4) \\
            &h(x)=(ax-4)^2
        \end{align*}
    \endgroup
    とおくとき，以下の問いに答えよ．

    \begin{enumerate}
        \item 放物線 $y=f(x)$ の頂点の座標は%
            $\left(\:\text{\myboxd{　ア　}}\,,\,\;\text{\myboxd{　イ　}}\:\right)$である．%
            $y=f(x)$ のグラフを$x$軸方向に\:\myboxd{　ウ　}\:だけ平行移動すると，$y=h(x)$ の%
            グラフに重なる．
        \item 放物線 $y=g(x)$ の頂点の座標は\\
            $\left(\:\text{\myboxb{　エ　}}\,,\,\;\text{\myboxb{　オ　}}\:a+%
            \:\text{\myboxb{\hspace{.5\zw}カキ\hspace{.5\zw}}}\:\right)$である．
        \item $y=f(x)$と$y=g(x)$のグラフの共有点はただ1つであり，その$x$座標は%
            \:\myboxd{　ク　}\:である．$y=f(x)$ と $y=h(x)$ のグラフの共有点はただ1つであり，%
            その座標は%
            $\left(\:\text{\myboxb{　ケ　}}\,,%
            \,\;\text{\myboxb{\hspace{.5\zw}コサ\hspace{.5\zw}}}\:\right)$である．%

            \vspace{.5\baselineskip}
            \textsf{ア，イ，ウ，クの解答群}
            
            　\begin{tabularx}{\dimexpr\linewidth-2\zw}{@{}LLLL@{}}
                ⓪\; $0$ & ①\; $1$ & ②\; $-1$ & ③\; $4$ \\ [4pt]
                ④\; $-4$ & ⑤\; $\myfracb{2}{a}$ &%
                ⑥\; $-\myfracb{2}{a}$ & ⑦\; $\myfracb{4}{a}$ \\ [9pt]
                ⑧\; $-\myfracb{4}{a}$ & ⑨\; $\myfracb{8}{a}$
            \end{tabularx}
        \item $y=f(x)$と$y=h(x)$のグラフと$x$軸に囲まれた図形を$A$とおく．%
            $A$の面積は $\myfraca{\text{\myboxb{シスセ}}}{\text{{\myboxb{　ソ　}}}\:a}$ である．%
            $y=f(x)$と$y=g(x)$と$y=h(x)$のグラフに囲まれた図形を$B$とおく．%
            $B$の面積は $\myfraca{\text{\myboxb{　タ　}}\:a}{\text{{\myboxb{　チ　}}}}$ である．%
            $A$の面積と$B$の面積が等しいときの定数$a$の値は\:\myboxb{　ツ　}\:である．
    \end{enumerate}
\end{document}
